\section{Especificaciones de casos de uso}

\subsection{Controles de simulación}

Los siguientes casos de uso permiten al usuario controlar la evolución del autómata celular. Esto implica evolucionar el sistema hacia adelante, hacia atrás y definiendo velocidades.

\subsubsection*{UC01 Iniciar simulación}

\begin{table}[H]
\begin{center}
	\begin{tabular}{| p{0.2\textwidth} | p{0.7\textwidth} |}
		\hline 
			\textbf{Actor} & Usuario. \\ \hline
			\textbf{Descripción} & El sistema inicia la evolución en pasos de tiempo constante. \\ \hline
			\textbf{Tipo} & Primario. \\ \hline
			\textbf{Entradas} &
			\begin{itemize}[noitemsep, nosep]
			\item Desde mouse \BTN{iniciar_simulacion}.
				\end{itemize}
			\\ \hline
			\textbf{Salidas} & Ninguna. \\ \hline
			\textbf{Precondición} & \begin{itemize}[noitemsep, nosep]
				\item Simulación detenida.
				\item Botón \BTN{iniciar_simulacion} habilitado.
			\end{itemize}
			\\ \hline
			\textbf{Postcondición} & \begin{itemize}[noitemsep, nosep]
				\item Simulación en ejecución.
				\item Botón \BTN{iniciar_simulacion} deshabilitado.
				\item Botón \BTN{pausar_simulacion} habilitado.
			\end{itemize} \\ \hline
			\textbf{Errores}  & Ninguno.  \\ \hline
			\end{tabular}
		\end{center}
	\end{table}

\begin{table}[H]
\begin{center}
	\begin{tabular}{|p{0.45\textwidth}|p{0.45\textwidth}|}
		\hline
			\multicolumn{2}{|c|}{\textbf{Flujo básico}} \\ \hline
			\textbf{Actor} & \textbf{Sistema} \\
			\hline
			\begin{enumerate}[noitemsep, nosep]
			\item Presiona el botón \BTN{iniciar_simulacion} desde \UI{control_ui}.
				\end{enumerate} 
			& 
			\begin{enumerate}[noitemsep, nosep]
				\item[2.] Deshabilita el botón \BTN{iniciar_simulacion}.
				\item[3.] Habilita el botón \BTN{pausar_simulacion}.
				\item[4.] Inicia la simulación. Cada paso evolutivo se realiza a una velocidad constante. 			
				\item[5.] Precarga evoluciones futuras en una cola.
				\item[6.] Actualiza de forma constante \UI{simulacion_ui} con la evolución del sistema.
				\end{enumerate} \\
			\hline
			\end{tabular} \\
		\end{center}
	\end{table}

\subsubsection*{UC02 Pausar simulación}

\begin{table}[H]
\begin{center}
	\begin{tabular}{| p{0.2\textwidth} | p{0.7\textwidth} |}
		\hline 
			\textbf{Actor} & Usuario. \\ \hline
			\textbf{Descripción} & Se detiene la simulación.  \\ \hline
			\textbf{Tipo} & Secundario. \\ \hline
			\textbf{Entradas} &
			\begin{itemize}[noitemsep, nosep]
			\item Desde mouse \BTN{pausar_simulacion}.
				\end{itemize}
			\\ \hline
			\textbf{Salidas} & Ninguna. \\ \hline
			\textbf{Precondición} &  \begin{itemize}[noitemsep, nosep]
				\item Simulación en ejecución.
				\item Botón \BTN{pausar_simulacion} habilitado.
			\end{itemize}			\\ \hline
			\textbf{Postcondición} & \begin{itemize}[noitemsep, nosep]
				\item Simulación detenida.
				\item Botón \BTN{pausar_simulacion} deshabilitado.
				\item Botón \BTN{iniciar_simulacion} habilitado.
			\end{itemize}  \\ \hline
			\textbf{Errores} & Ninguno. \\ \hline
			\end{tabular}
			\end{center}
				\end{table}
			
		\begin{table}[H]
	\begin{center}
\begin{tabular}{|p{0.45\textwidth}|p{0.45\textwidth}|}
\hline
\multicolumn{2}{|c|}{\textbf{Flujo básico}} \\ \hline
	\textbf{Actor} & \textbf{Sistema} \\
		\hline
			\begin{enumerate}[noitemsep, nosep]
			\item Presiona el botón \BTN{pausar_simulacion} desde \UI{control_ui}.
			\end{enumerate} 
			& 
			\begin{enumerate}[noitemsep, nosep]
				\item[2.] Detiene la simulación en \UI{simulacion_ui}. 
				\item[3.] Deshabilita el botón \BTN{pausar_simulacion}.
				\item[4.] Habilita el botón \BTN{iniciar_simulacion}. 
			\end{enumerate} \\
				\hline
			\end{tabular} \\
			\end{center}
		\end{table}
	
\subsubsection*{UC03 Cambiar velocidad}

\begin{table}[H]
	\begin{center}
		\begin{tabular}{| p{0.2\textwidth} | p{0.7\textwidth} |}
			\hline 
			\textbf{Actor} & Usuario. \\ \hline
			\textbf{Descripción} & Se cambia la velocidad de actualización del autómata celular. \\ \hline
			\textbf{Tipo} & Secundario. \\ \hline
			\textbf{Entradas} &
			\begin{itemize}[noitemsep, nosep]
				\item Desde mouse \SLD{velocidad}.
			\end{itemize}
			\\ \hline
			\textbf{Salidas} & Ninguna. \\ \hline
			\textbf{Precondición} & 
				Ninguna.
			\\ \hline
			\textbf{Postcondición} & \begin{itemize}[noitemsep, nosep]
				\item Cambia la velocidad de simulación.
			\end{itemize}  \\ \hline
			\textbf{Errores} & 
				Ninguno. \\ \hline
			\end{tabular}
	\end{center}
\end{table}

\begin{table}[H]
	\begin{center}
		\begin{tabular}{|p{0.45\textwidth}|p{0.45\textwidth}|}
			\hline
			\multicolumn{2}{|c|}{\textbf{Flujo básico}} \\ \hline
			\textbf{Actor} & \textbf{Sistema} \\
			\hline 
			\begin{enumerate}[noitemsep, nosep]
				\item Presiona \SLD{velocidad} desde \UI{control_ui} y desliza a la velocidad deseada.
			\end{enumerate} 
			& 
			\begin{enumerate}[noitemsep, nosep] 
				\item[2.] Cambia la velocidad de la simulación a la velocidad seleccionada.
			\end{enumerate} \\
			\hline
		\end{tabular} \\
	\end{center}
\end{table}

\subsubsection*{UC04 Simular un paso}

\begin{table}[H]
\begin{center}
\begin{tabular}{| p{0.2\textwidth} | p{0.7\textwidth} |}
	\hline 
		\textbf{Actor} & Usuario. \\ \hline
			\textbf{Descripción} & Se realiza la simulación del CA, pero únicamente en un paso. \\ \hline
			\textbf{Tipo} & Primario. \\ \hline
			\textbf{Entradas} &
			\begin{itemize}[noitemsep, nosep]
			\item Desde mouse \BTN{paso_adelante}.
			\end{itemize}
				\\ \hline
			\textbf{Salidas} & Ninguna. \\ \hline
			\textbf{Precondición} & Ninguna.
			\\ \hline
			\textbf{Postcondición} & \begin{itemize}[noitemsep, nosep]
				\item Evoluciona el sistema.
			\end{itemize}  \\ \hline
			\textbf{Errores} & 
			\begin{itemize}[noitemsep, nosep]
			\item Velocidad máxima: Se intenta dar un paso cuando la simulación está activa en la velocidad máxima.
			\end{itemize}\\ \hline
				\end{tabular}
			\end{center}
		\end{table}
	
\begin{table}[H]
\begin{center}
\begin{tabular}{|p{0.45\textwidth}|p{0.45\textwidth}|}
	\hline
		\multicolumn{2}{|c|}{\textbf{Flujo básico}} \\ \hline
			\textbf{Actor} & \textbf{Sistema} \\
			\hline
			\begin{enumerate}[noitemsep, nosep]
			\item Presiona el botón \BTN{paso_adelante} desde \UI{control_ui}.
			\end{enumerate} 
				& 
			\begin{enumerate}[noitemsep, nosep]
			\item[2.] La simulación avanza un paso.
			\item[3.] Actualiza \UI{simulacion_ui}.
			\end{enumerate} \\
				\hline
			\end{tabular} \\
			\end{center}
		\end{table}
	
\begin{table}[H]
\begin{center}
\begin{tabular}{|p{0.45\textwidth}|p{0.45\textwidth}|}
	\hline
		\multicolumn{2}{|c|}{\textbf{Flujo alterno 1: Velocidad máxima}} \\ \hline
			\textbf{Actor} & \textbf{Sistema} \\
			\hline
			\begin{enumerate}[noitemsep, nosep]
			\item[1.1.] Presiona el botón \BTN{paso_adelante} desde \UI{control_ui} cuando la simulación está activa y la velocidad está en la velocidad máxima.
			\end{enumerate} 
				& 
			\begin{enumerate}[noitemsep, nosep]
			\item[1.2.] Agrega la simulación a una cola y mantiene la velocidad actual. 
			\item[1.3.] Actualiza \UI{simulacion_ui}.
			\end{enumerate} \\
				\hline
			\end{tabular} 
			\end{center}
		\end{table}
		
\subsubsection*{UC05 Retroceder un paso}

\begin{table}[H]
	\begin{center}
		\begin{tabular}{| p{0.2\textwidth} | p{0.7\textwidth} |}
			\hline 
			\textbf{Actor} & Usuario. \\ \hline
			\textbf{Descripción} & El sistema restaura el estado anterior de la simulación. \\ \hline
			\textbf{Tipo} & Secundario. \\ \hline
			\textbf{Entradas} &
			\begin{itemize}[noitemsep, nosep]
				\item Desde mouse \BTN{paso_atras}.
			\end{itemize}
			\\ \hline
			\textbf{Salidas} & Ninguna. \\ \hline
			\textbf{Precondición} & \begin{itemize}[noitemsep, nosep]
				\item Existe un estado previo.
			\end{itemize}
			\\ \hline
			\textbf{Postcondición} & \begin{itemize}[noitemsep, nosep]
				\item El sistema restaura la evolución pasada.
			\end{itemize}  \\ \hline
			\textbf{Errores} & 
			\begin{itemize}[noitemsep, nosep]
				\item Configuración inicial: El usuario intenta retroceder la simulación de la configuración inicial.
			\end{itemize}\\ \hline
		\end{tabular}
	\end{center}
\end{table}

\begin{table}[H]
	\begin{center}
		\begin{tabular}{|p{0.45\textwidth}|p{0.45\textwidth}|}
			\hline
			\multicolumn{2}{|c|}{\textbf{Flujo básico}} \\ \hline
			\textbf{Actor} & \textbf{Sistema} \\
			\hline
			\begin{enumerate}[noitemsep, nosep]
				\item Presiona el botón \BTN{paso_atras} desde \UI{control_ui}.
			\end{enumerate} 
			& 
			\begin{enumerate}[noitemsep, nosep]
				\item[2.] La simulación retrocede un paso. 
				\item[3.] Actualiza \UI{simulacion_ui}. 
			\end{enumerate} \\
			\hline
		\end{tabular} \\
	\end{center}
\end{table}

\begin{table}[H]
	\begin{center}
		\begin{tabular}{|p{0.45\textwidth}|p{0.45\textwidth}|}
			\hline
			\multicolumn{2}{|c|}{\textbf{Flujo alterno 1: Configuración inicial}} \\ \hline
			\textbf{Actor} & \textbf{Sistema} \\
			\hline
			\begin{enumerate}[noitemsep, nosep]
				\item[1.1.] Presiona el botón \BTN{paso_atras} desde \UI{control_ui} estando en la configuración inicial (no hay más configuraciones antes).  
			\end{enumerate}
			& 
			\begin{enumerate}[noitemsep, nosep]
				\item[1.2.] No realiza ningún cambio.
			\end{enumerate} \\
			\hline
		\end{tabular} 
	\end{center}
\end{table}

\subsubsection*{UC06 Reiniciar simulación}

\begin{table}[H]
	\begin{center}
		\begin{tabular}{| p{0.2\textwidth} | p{0.7\textwidth} |}
			\hline 
			\textbf{Actor} & Usuario. \\ \hline
			\textbf{Descripción} & Restablece la simulación a su estado inicial (todas las células vacías). \\ \hline
			\textbf{Tipo} & Secundario. \\ \hline
			\textbf{Entradas} &
			\begin{itemize}[noitemsep, nosep]
				\item Desde mouse \BTN{reiniciar_simulacion}.
			\end{itemize}
			\\ \hline
			\textbf{Salidas} & Ninguna. \\ \hline
			\textbf{Precondición} & Ninguna.
			\\ \hline
			\textbf{Postcondición} & \begin{itemize}[noitemsep, nosep]
				\item El sistema adquiere la configuración nula.
			\end{itemize}  \\ \hline
			\textbf{Errores} & Ninguno. \\ \hline
		\end{tabular}
	\end{center}
\end{table}

\begin{table}[H]
	\begin{center}
		\begin{tabular}{|p{0.45\textwidth}|p{0.45\textwidth}|}
			\hline
			\multicolumn{2}{|c|}{\textbf{Flujo básico}} \\ \hline
			\textbf{Actor} & \textbf{Sistema} \\
			\hline
			\begin{enumerate}[noitemsep, nosep]
				\item Presiona el botón \BTN{reiniciar_simulacion} desde \UI{control_ui}.
			\end{enumerate} 
			& 
			\begin{enumerate}[noitemsep, nosep]
				\item[2.] Regresa la configuración a la última versión modificada mientras estaba en pausa. 
				\item[3.] Elimina el historial de evoluciones.
				\item[4.] Actualiza \UI{simulacion_ui}.
			\end{enumerate} \\
			\hline
		\end{tabular} \\
	\end{center}
\end{table}

\subsubsection*{UC07 Aleatorizar espacio}

\begin{table}[H]
	\begin{center}
		\begin{tabular}{| p{0.2\textwidth} | p{0.7\textwidth} |}
			\hline 
			\textbf{Actor} & Usuario. \\ \hline
			\textbf{Descripción} & Establece una configuración aleatoria en el espacio. \\ \hline
			\textbf{Tipo} & Secundario. \\ \hline
			\textbf{Entradas} &
			\begin{itemize}[noitemsep, nosep]
				\item Desde mouse \BTN{aleatorizar_simulacion}.
			\end{itemize}
			\\ \hline
			\textbf{Salidas} & Ninguna. \\ \hline
			\textbf{Precondición} & Ninguna.
			\\ \hline
			\textbf{Postcondición} & \begin{itemize}[noitemsep, nosep]
				\item El sistema adquiere una configuración aleatoria.
			\end{itemize}  \\ \hline
			\textbf{Errores} & Ninguno. \\ \hline
		\end{tabular}
	\end{center}
\end{table}

\begin{table}[H]
	\begin{center}
		\begin{tabular}{|p{0.45\textwidth}|p{0.45\textwidth}|}
			\hline
			\multicolumn{2}{|c|}{\textbf{Flujo básico}} \\ \hline
			\textbf{Actor} & \textbf{Sistema} \\
			\hline
			\begin{enumerate}[noitemsep, nosep]
				\item Presiona el botón \BTN{aleatorizar_simulacion} desde \UI{control_ui}.
			\end{enumerate} 
			& 
			\begin{enumerate}[noitemsep, nosep]
				\item[2.] La configuración del espacio celular adquiere valores aleatorios según el número de estados definidos.
				\item[3.] Elimina el historial de evoluciones.
				\item[4.] Actualiza \UI{simulacion_ui}.
			\end{enumerate} \\
			\hline
		\end{tabular} \\
	\end{center}
\end{table}

\subsubsection*{UC08 Importar evolución}

\begin{table}[H]
	\begin{center}
		\begin{tabular}{| p{0.2\textwidth} | p{0.7\textwidth} |}
			\hline 
			\textbf{Actor} & Usuario. \\ \hline
			\textbf{Descripción} & Importa una evolución guardada en formato ``npy''. \\ \hline
			\textbf{Tipo} & Secundario. \\ \hline
			\textbf{Entradas} &
			\begin{itemize}[noitemsep, nosep]
				\item Desde mouse \BTN{importar}.
			\end{itemize}
			\\ \hline
			\textbf{Salidas} & \begin{itemize}[noitemsep, nosep]
				\item \MSG{formato_incorrecto}.
				\item \MSG{tamaño_incorrecto}.
			\end{itemize} \\ \hline
			\textbf{Precondición} & Ninguna.
			\\ \hline
			\textbf{Postcondición} & \begin{itemize}[noitemsep, nosep]
				\item El sistema carga la simulación y la muestra en \UI{simulacion_ui}.
			\end{itemize}  \\ \hline
			\textbf{Errores} & \begin{itemize}[noitemsep, nosep]
				\item Formato inválido: El archivo no coincide.
				\item Tamaño no compatible: El tamaño no coincide.
			\end{itemize}  \\ \hline
		\end{tabular}
	\end{center}
\end{table}

\begin{table}[H]
	\begin{center}
		\begin{tabular}{|p{0.45\textwidth}|p{0.45\textwidth}|}
			\hline
			\multicolumn{2}{|c|}{\textbf{Flujo básico}} \\ \hline
			\textbf{Actor} & \textbf{Sistema} \\
			\hline
			\begin{enumerate}[noitemsep, nosep]
				\item Presiona el botón \BTN{importar} desde \UI{control_ui}.
				\item[3. ] Selecciona el archivo ``npy'' que contenga una evolución.
			\end{enumerate} 
			& 
			\begin{enumerate}[noitemsep, nosep]
				\item[2.] Abre el explorador de archivos.
				\item[4.] Carga la evolución.
				\item[5.] Actualiza \UI{simulacion_ui}.
			\end{enumerate} \\
			\hline
		\end{tabular} \\
	\end{center}
\end{table}

\begin{table}[H]
	\begin{center}
		\begin{tabular}{|p{0.45\textwidth}|p{0.45\textwidth}|}
			\hline
			\multicolumn{2}{|c|}{\textbf{Flujo alterno 1: Formato inválido}} \\ \hline
			\textbf{Actor} & \textbf{Sistema} \\
			\hline
			\begin{enumerate}[noitemsep, nosep]
				\item[1.1.] Presiona el botón \BTN{importar} desde \UI{control_ui}.
				\item[1.3. ]   Selecciona un archivo ``npy'' inválido u otro tipo de archivo.
			\end{enumerate}
			& 
			\begin{enumerate}[noitemsep, nosep]
				\item[1.2.] Abre el explorador de archivos.
				\item[1.4.] No realiza ningún cambio y muestra \MSG{formato_incorrecto}.
			\end{enumerate} \\
			\hline
		\end{tabular} 
	\end{center}
\end{table}

\begin{table}[H]
	\begin{center}
		\begin{tabular}{|p{0.45\textwidth}|p{0.45\textwidth}|}
			\hline
			\multicolumn{2}{|c|}{\textbf{Flujo alterno 2: Tamaño no compatible}} \\ \hline
			\textbf{Actor} & \textbf{Sistema} \\
			\hline
			\begin{enumerate}[noitemsep, nosep]
				\item[2.1.] Presiona el botón \BTN{importar} desde \UI{control_ui}.
				\item[2.3. ] Selecciona el archivo ``npy'' que contenga una evolución con un tamaño diferente a \UI{simulacion_ui}.
			\end{enumerate}
			& 
			\begin{enumerate}[noitemsep, nosep]
				\item[2.2.] Abre el explorador de archivos.
				\item[2.4.] Recorta o centra la evolución y muestra el mensaje \MSG{tamaño_incorrecto}.
				\item[2.5.] Actualiza \UI{simulacion_ui}.
			\end{enumerate} \\
			\hline
		\end{tabular} 
	\end{center}
\end{table}

\subsubsection*{UC09 Grabar evolución}

\begin{table}[H]
	\begin{center}
		\begin{tabular}{| p{0.2\textwidth} | p{0.7\textwidth} |}
			\hline 
			\textbf{Actor} & Usuario. \\ \hline
			\textbf{Descripción} & Evoluciona el sistema sin visualización y guarda todos los pasos. \\ \hline
			\textbf{Tipo} & Secundario. \\ \hline
			\textbf{Entradas} &
			\begin{itemize}[noitemsep, nosep]
				\item Desde mouse \BTN{grabar_evolucion}.
				\item Desde teclado \IN{grabar_pasos}.
			\end{itemize}
			\\ \hline
			\textbf{Salidas} & Ninguna. \\ \hline
			\textbf{Precondición} & Ninguna.
			\\ \hline
			\textbf{Postcondición} & \begin{itemize}[noitemsep, nosep]
				\item El sistema evoluciona los pasos seleccionados.
			\end{itemize}  \\ \hline
			\textbf{Errores} & Ninguno. \\ \hline
		\end{tabular}
	\end{center}
\end{table}

\begin{table}[H]
	\begin{center}
		\begin{tabular}{|p{0.45\textwidth}|p{0.45\textwidth}|}
			\hline
			\multicolumn{2}{|c|}{\textbf{Flujo básico}} \\ \hline
			\textbf{Actor} & \textbf{Sistema} \\
			\hline
			\begin{enumerate}[noitemsep, nosep]
				\item Presiona el botón \BTN{grabar_evolucion} desde \UI{control_ui}.
				\item[3. ] Elige un número con \IN{grabar_pasos}.
			\end{enumerate} 
			& 
			\begin{enumerate}[noitemsep, nosep]
				\item[2.] Abre una ventana preguntando el número de pasos a evolucionar.
				\item[4.] Evoluciona el sistema el número de pasos.
				\item[5.] Muestra en \UI{simulacion_ui} el último paso.
			\end{enumerate} \\
			\hline
		\end{tabular} \\
	\end{center}
\end{table}

\subsubsection*{UC10 Exportar evolución}

\begin{table}[H]
	\begin{center}
		\begin{tabular}{| p{0.2\textwidth} | p{0.7\textwidth} |}
			\hline 
			\textbf{Actor} & Usuario. \\ \hline
			\textbf{Descripción} & Exporta la evolución en formato ``npy'' y ``mat''. \\ \hline
			\textbf{Tipo} & Primario. \\ \hline
			\textbf{Entradas} &
			\begin{itemize}[noitemsep, nosep]
				\item Desde mouse \BTN{exportar}.
			\end{itemize}
			\\ \hline
			\textbf{Salidas} & \begin{itemize}[noitemsep, nosep]
				\item Evolución en formato ``npy''.
				\item Evolución en formato ``mat''.
			\end{itemize} \\ \hline
			\textbf{Precondición} & UC09 Grabar evolución.
			\\ \hline
			\textbf{Postcondición} & Ninguna.  \\ \hline
			\textbf{Errores} & \begin{itemize}[noitemsep, nosep]
				\item Sin evolución: No se ha grabado ninguna evolución.
			\end{itemize}  \\ \hline
		\end{tabular}
	\end{center}
\end{table}

\begin{table}[H]
	\begin{center}
		\begin{tabular}{|p{0.45\textwidth}|p{0.45\textwidth}|}
			\hline
			\multicolumn{2}{|c|}{\textbf{Flujo básico}} \\ \hline
			\textbf{Actor} & \textbf{Sistema} \\
			\hline
			\begin{enumerate}[noitemsep, nosep]
				\item Presiona el botón \BTN{exportar} desde \UI{control_ui}.
				\item[3. ] Selecciona la ruta de guardado archivo.
			\end{enumerate} 
			& 
			\begin{enumerate}[noitemsep, nosep]
				\item[2.] Abre el explorador de archivos.
				\item[4.] Crea los archivos de la evolución en la ruta especificada.
			\end{enumerate} \\
			\hline
		\end{tabular} \\
	\end{center}
\end{table}

\begin{table}[H]
	\begin{center}
		\begin{tabular}{|p{0.45\textwidth}|p{0.45\textwidth}|}
			\hline
			\multicolumn{2}{|c|}{\textbf{Flujo alterno 1: Sin evolución}} \\ \hline
			\textbf{Actor} & \textbf{Sistema} \\
			\hline
			\begin{enumerate}[noitemsep, nosep]
				\item[1.1.] Presiona el botón \BTN{exportar} desde \UI{control_ui} sin haber una evolución.
			\end{enumerate}
			& 
			\begin{enumerate}[noitemsep, nosep]
				\item[1.2.] No realiza ningún cambio y muestra \MSG{evolucion_vacia}.
			\end{enumerate} \\
			\hline
		\end{tabular} 
	\end{center}
\end{table}

\subsubsection*{UC11 Cargar pesos}

\begin{table}[H]
	\begin{center}
		\begin{tabular}{| p{0.2\textwidth} | p{0.7\textwidth} |}
			\hline 
			\textbf{Actor} & Usuario. \\ \hline
			\textbf{Descripción} & Carga los pesos de la red neuronal entrenada en el autómata celular neuronal. \\ \hline
			\textbf{Tipo} & Primario. \\ \hline
			\textbf{Entradas} &
			\begin{itemize}[noitemsep, nosep]
				\item Desde mouse \BTN{cargar_pesos}.
			\end{itemize}
			\\ \hline
			\textbf{Salidas} & Ninguna. \\ \hline
			\textbf{Precondición} & Ninguna.
			\\ \hline
			\textbf{Postcondición} & \begin{itemize}[noitemsep, nosep]
				\item Habilita el botón \BTN{cambiar_ca_nca}.
			\end{itemize}  \\ \hline
			\textbf{Errores} & \begin{itemize}[noitemsep, nosep]
				\item Formato inválido: El archivo no coincide.
			\end{itemize}  \\ \hline
		\end{tabular}
	\end{center}
\end{table}

\begin{table}[H]
	\begin{center}
		\begin{tabular}{|p{0.45\textwidth}|p{0.45\textwidth}|}
			\hline
			\multicolumn{2}{|c|}{\textbf{Flujo básico}} \\ \hline
			\textbf{Actor} & \textbf{Sistema} \\
			\hline
			\begin{enumerate}[noitemsep, nosep]
				\item Presiona el botón \BTN{cargar_pesos} desde \UI{control_ui}.
				\item[3. ] Selecciona los pesos del archivo ``npz''.
			\end{enumerate} 
			& 
			\begin{enumerate}[noitemsep, nosep]
				\item[2.] Abre el explorador de archivos.
				\item[4.] Carga los pesos del autómata celular neuronal.
				\item[5.] Habilita el botón \BTN{cambiar_ca_nca}.
			\end{enumerate} \\
			\hline
		\end{tabular} \\
	\end{center}
\end{table}

\begin{table}[H]
	\begin{center}
		\begin{tabular}{|p{0.45\textwidth}|p{0.45\textwidth}|}
			\hline
			\multicolumn{2}{|c|}{\textbf{Flujo alterno 1: Formato inválido}} \\ \hline
			\textbf{Actor} & \textbf{Sistema} \\
			\hline
			\begin{enumerate}[noitemsep, nosep]
				\item[1.1.] Presiona el botón \BTN{cargar_pesos} desde \UI{control_ui}.
				\item[1.3.]  Selecciona un archivo ``npz'' inválido u otro tipo de archivo.
			\end{enumerate}
			& 
			\begin{enumerate}[noitemsep, nosep]
				\item[1.2.] Abre el explorador de archivos.
				\item[1.4.] No realiza ningún cambio y muestra \MSG{formato_incorrecto}.
			\end{enumerate} \\
			\hline
		\end{tabular} 
	\end{center}
\end{table}

\subsubsection*{UC12 Cambiar CA/NCA}


\begin{table}[H]
	\begin{center}
		\begin{tabular}{| p{0.2\textwidth} | p{0.7\textwidth} |}
			\hline 
			\textbf{Actor} & Usuario. \\ \hline
			\textbf{Descripción} & Cambia la evolución del sistema entre el autómata celular y el autómata celular neuronal. \\ \hline
			\textbf{Tipo} & Secundario. \\ \hline
			\textbf{Entradas} &
			\begin{itemize}[noitemsep, nosep]
				\item Desde mouse \BTN{cambiar_ca_nca}.
			\end{itemize}
			\\ \hline
			\textbf{Salidas} & Ninguna. \\ \hline
			\textbf{Precondición} & UC11 Cargar pesos.
			\\ \hline
			\textbf{Postcondición} & Ninguna.  \\ \hline
			\textbf{Errores} & Ninguno.  \\ \hline
		\end{tabular}
	\end{center}
\end{table}

\begin{table}[H]
	\begin{center}
		\begin{tabular}{|p{0.45\textwidth}|p{0.45\textwidth}|}
			\hline
			\multicolumn{2}{|c|}{\textbf{Flujo básico}} \\ \hline
			\textbf{Actor} & \textbf{Sistema} \\
			\hline
			\begin{enumerate}[noitemsep, nosep]
				\item Presiona el botón \BTN{cambiar_ca_nca} desde \UI{control_ui}.
			\end{enumerate} 
			& 
			\begin{enumerate}[noitemsep, nosep]
				\item[2.] Cambia entre autómata celular y autómata celular neuronal.
			\end{enumerate} \\
			\hline
		\end{tabular} \\
	\end{center}
\end{table}

\subsection{Configuración del autómata celular}

Estos casos de uso definen las características generales del autómata celular, tales como el número de estados, vecindad, comportamiento de frontera y la función de transición.

\subsubsection*{UC13 Definir vecindad}

	\begin{table}[H]
		\begin{center}
			\begin{tabular}{| p{0.2\textwidth} | p{0.7\textwidth} |}
				\hline 
				\textbf{Actor} & Usuario. \\ \hline
				\textbf{Descripción} & Se define el conjunto de posiciones relativas de cada célula que forman parte de la vecindad.  \\ \hline
				\textbf{Tipo} & Primario. \\ \hline
				\textbf{Entradas} &
				\begin{itemize}[noitemsep, nosep]
					\item Desde mouse \BTN{def_vecindad}.
					\item Desde mouse \BTN{añadir_grupo}.
					\item Desde mouse \BTN{eliminar_grupo}.
					\item Desde mouse \IN{grupo_vecindad}.
					\item Desde mouse \IN{espacio_vecindad}.
				\end{itemize}
				\\ \hline
				\textbf{Salidas} & \begin{itemize}[noitemsep, nosep]
						\item Se actualiza \IN{regla_codigo}.
					\end{itemize} \\ \hline
				\textbf{Precondición} &  Ninguna.
				\\ \hline
				\textbf{Postcondición} &  \begin{itemize}[noitemsep, nosep]
					\item Vecindad actualizada en el sistema.
					\item Simulación en pausa.
				\end{itemize} \\ \hline
				\textbf{Errores} & \begin{itemize}[noitemsep, nosep]
					\item Vecindad vacía: No se seleccionó ningún elemento en \IN{espacio_vecindad}.
				\end{itemize} \\ \hline
			\end{tabular}
		\end{center}
	\end{table}
		
	\begin{table}[H]
		\begin{center}
			\begin{tabular}{|p{0.45\textwidth}|p{0.45\textwidth}|}
				\hline
				\multicolumn{2}{|c|}{\textbf{Flujo básico}} \\ \hline
				\textbf{Actor} & \textbf{Sistema} \\
				\hline
				\begin{enumerate}[noitemsep, nosep]
					\item Presiona el \IN{grupo_vecindad} en \UI{configuracion_ui}. 
					\item Selecciona múltiples elementos del \IN{espacio_vecindad}.
					\item[4.] Presiona el botón \BTN{def_vecindad}.
				\end{enumerate} 
				& 
				\begin{enumerate}[noitemsep, nosep]
					\item[3.] Asigna el grupo seleccionado al elemento de \IN{espacio_vecindad}.
					\item[5.] Detiene la simulación.
					\item[6.] Asigna la nueva vecindad y actualiza el tamaño de reglas permitidas.
					\item[7.] Actualiza \IN{regla_codigo}.
				\end{enumerate} \\
				\hline
			\end{tabular} 
		\end{center}
	\end{table}
	
	\begin{table}[H]
		\begin{center}
			\begin{tabular}{|p{0.45\textwidth}|p{0.45\textwidth}|}
				\hline
				\multicolumn{2}{|c|}{\textbf{Flujo alterno 1: Vecindad vacía}} \\ \hline
				\textbf{Actor} & \textbf{Sistema} \\
				\hline
				\begin{enumerate}[noitemsep, nosep] 
					\item[1.2.] Selecciona el grupo vacío.
					\item[1.4.] Presiona el botón \BTN{def_vecindad}.
				\end{enumerate} 
				& 
				\begin{enumerate}[noitemsep, nosep]
					\item[1.3.] Asigna el grupo a todos los elementos de \IN{espacio_vecindad}.
					\item[1.5.] No detiene la animación ni actualiza la vecindad.
				\end{enumerate} \\
				\hline
			\end{tabular} 
		\end{center}
	\end{table}
	
	\begin{table}[H]
		\begin{center}
			\begin{tabular}{|p{0.45\textwidth}|p{0.45\textwidth}|}
				\hline
				\multicolumn{2}{|c|}{\textbf{Flujo alterno 2: Añadir grupo}} \\ \hline
				\textbf{Actor} & \textbf{Sistema} \\
				\hline
				\begin{enumerate}[noitemsep, nosep]
					\item[2.1.] Presiona el botón \BTN{añadir_grupo}.
				\end{enumerate} 
				& 
				\begin{enumerate}[noitemsep, nosep]
					\item[2.2.] Añade un nuevo \IN{grupo_vecindad}.
				\end{enumerate} \\
				\hline
			\end{tabular} 
		\end{center}
	\end{table}

	\begin{table}[H]
		\begin{center}
			\begin{tabular}{|p{0.45\textwidth}|p{0.45\textwidth}|}
				\hline
				\multicolumn{2}{|c|}{\textbf{Flujo alterno 3: Eliminar grupo}} \\ \hline
				\textbf{Actor} & \textbf{Sistema} \\
				\hline
				\begin{enumerate}[noitemsep, nosep]
					\item[3.1.] Presiona el botón \BTN{eliminar_grupo}.
				\end{enumerate} 
				& 
				\begin{enumerate}[noitemsep, nosep]
					\item[3.2.] Borra todas las celdas en \IN{espacio_vecindad} asignadas al último \IN{grupo_vecindad}.
					\item[3.3.] Elimina el último \IN{grupo_vecindad}.
				\end{enumerate} \\
				\hline
			\end{tabular} 
		\end{center}
	\end{table}

\subsubsection*{UC14 Definir estados}

	\begin{table}[H]
		\begin{center}
			\begin{tabular}{| p{0.2\textwidth} | p{0.7\textwidth} |}
				\hline 
				\textbf{Actor} & Usuario. \\ \hline
				\textbf{Descripción} & Se define el número de estados.  \\ \hline
				\textbf{Tipo} & Secundario. \\ \hline
				\textbf{Entradas} &
				\begin{itemize}[noitemsep, nosep]
					\item Desde teclado \IN{estados}.
					\item Desde mouse \BTN{def_estados}.
				\end{itemize}
				\\ \hline
				\textbf{Salidas} &  \begin{itemize}[noitemsep, nosep]
					\item Se añaden \BTN{estado}.
					\item Se actualiza \IN{regla_codigo}.
				\end{itemize}\\ \hline
				\textbf{Precondición} & Ninguna. \\ \hline
				\textbf{Postcondición} & \begin{itemize}[noitemsep, nosep]
					\item Estados actualizados en el sistema.
					\item Simulación en pausa.
					\item Botones de estados añadidos.
				\end{itemize} \\ \hline
				\textbf{Errores} &  \begin{itemize}[noitemsep, nosep]
					\item Inválido: El número de estados ingresados es inválido.
				\end{itemize} \\ \hline
			\end{tabular}
		\end{center}
	\end{table}
	
	\begin{table}[H]
		\begin{center}
			\begin{tabular}{|p{0.45\textwidth}|p{0.45\textwidth}|}
				\hline
				\multicolumn{2}{|c|}{\textbf{Flujo básico}} \\ \hline
				\textbf{Actor} & \textbf{Sistema} \\
				\hline
				\begin{enumerate}[noitemsep, nosep]
					\item[1. ] Ingresa el número de estados en \IN{estados} en \UI{configuracion_ui} y presiona el botón \BTN{def_estados}.
				\end{enumerate} 
				& 
				\begin{enumerate}[noitemsep, nosep] 
					\item[2.] Se detiene la simulación.
					\item[3.] Se añaden \BTN{estado} definidos.
					\item[4.] Actualiza el tamaño de las reglas permitidas.
					\item[5.] Actualiza \IN{regla_codigo}.
				\end{enumerate} \\
				\hline
			\end{tabular} 
		\end{center}
	\end{table}
	
	\begin{table}[H]
		\begin{center}
			\begin{tabular}{|p{0.45\textwidth}|p{0.45\textwidth}|}
				\hline
				\multicolumn{2}{|c|}{\textbf{Flujo alterno 1: Inválido}} \\ \hline
				\textbf{Actor} & \textbf{Sistema} \\
				\hline
				\begin{enumerate}[noitemsep, nosep]
					\item[1.1. ] Ingresa el número de estados en \IN{estados} inválido.
				\end{enumerate} 
				& 
				\begin{enumerate}[noitemsep, nosep] 
					\item[1.2.] Corrige el número al valor más cercano posible.
				\end{enumerate} \\
				\hline
			\end{tabular} 
		\end{center}
	\end{table}

\subsubsection*{UC15 Definir condición de frontera}

\begin{table}[H]
	\begin{center}
		\begin{tabular}{| p{0.2\textwidth} | p{0.7\textwidth} |}
			\hline 
			\textbf{Actor} & Usuario. \\ \hline
			\textbf{Descripción} & Se selecciona entre diferentes opciones de condiciones de frontera. \\ \hline
			\textbf{Tipo} & Secundario. \\ \hline
			\textbf{Entradas} &
			\begin{itemize}[noitemsep, nosep]
				\item Desde mouse \BTN{def_frontera}.
			\end{itemize}
			\\ \hline
			\textbf{Salidas} & Ninguna.
			\\ \hline
			\textbf{Precondición} & Ninguna.
			\\ \hline
			\textbf{Postcondición} & \begin{itemize}[noitemsep, nosep]
				\item Condición actualizada en el sistema.
			\end{itemize} \\ \hline
			\textbf{Errores} & Ninguno. \\ \hline
		\end{tabular}
	\end{center}
\end{table}

\begin{table}[H]
	\begin{center}
		\begin{tabular}{|p{0.45\textwidth}|p{0.45\textwidth}|}
			\hline
			\multicolumn{2}{|c|}{\textbf{Flujo básico}} \\ \hline
			\textbf{Actor} & \textbf{Sistema} \\
			\hline
			\begin{enumerate}[noitemsep, nosep]
				\item Presiona el botón \BTN{def_frontera} desde \UI{configuracion_ui}.
			\end{enumerate} 
			& 
			\begin{enumerate}[noitemsep, nosep]
				\item[2.] Actualiza la condición de frontera a la siguiente configuración.  
			\end{enumerate} \\
			\hline
		\end{tabular} \\
	\end{center}
\end{table}

\subsubsection*{UC16 Definir regla de transición}
	
	\begin{table}[H]
		\begin{center}
			\begin{tabular}{| p{0.2\textwidth} | p{0.7\textwidth} |}
				\hline 
				\textbf{Actor} & Usuario. \\ \hline
				\textbf{Descripción} & Se configura la regla que define el comportamiento del CA. \\ \hline
				\textbf{Tipo} & Primario. \\ \hline
				\textbf{Entradas} &
				\begin{itemize}[noitemsep, nosep]
					\item Desde mouse \BTN{def_regla}.
					\item Desde teclado \IN{regla}.
					\item Desde mouse \IN{regla_codigo}.
				\end{itemize}
				\\ \hline
				\textbf{Salidas} &  \begin{itemize}[noitemsep, nosep]
					\item Se actualiza \IN{regla_codigo}.
					\item Se actualiza \IN{regla}.
				\end{itemize}     \\ \hline
				\textbf{Precondición} & Ninguna.
				\\ \hline
				\textbf{Postcondición} &  \begin{itemize}[noitemsep, nosep]
					\item Regla actualizada en el sistema.
					\item Simulación en pausa.
				\end{itemize}   \\ \hline
				\textbf{Errores} & 
				\begin{itemize}[noitemsep, nosep]
					\item Inválido: La regla establecida está fuera del rango. 
				\end{itemize}\\ \hline
			\end{tabular}
		\end{center}
	\end{table}
	
	\begin{table}[H]
		\begin{center}
			\begin{tabular}{|p{0.45\textwidth}|p{0.45\textwidth}|}
				\hline
				\multicolumn{2}{|c|}{\textbf{Flujo básico}} \\ \hline
				\textbf{Actor} & \textbf{Sistema} \\
				\hline
				\begin{enumerate}[noitemsep, nosep]
					\item Ingresa la regla en el campo \IN{regla} en \UI{configuracion_ui}.
					\item Presiona el botón \BTN{def_regla}.
				\end{enumerate} 
				& 
				\begin{enumerate}[noitemsep, nosep]
					\item[3.] Actualiza la regla.
					\item[4.] Actualiza \IN{regla_codigo}.
				\end{enumerate} \\
				\hline
			\end{tabular} 
				\end{center}
		\end{table}
		
		\begin{table}[H]
			\begin{center}
				\begin{tabular}{|p{0.45\textwidth}|p{0.45\textwidth}|}
					\hline
					\multicolumn{2}{|c|}{\textbf{Flujo alterno 1}} \\ \hline
					\textbf{Actor} & \textbf{Sistema} \\
					\hline
					\begin{enumerate}[noitemsep, nosep]
						\item[1.1.] Ingresa la regla seleccionando \IN{regla_codigo} en \UI{configuracion_ui}.
						\item[1.2.] Presiona el botón \BTN{def_regla}.
					\end{enumerate} 
					& 
					\begin{enumerate}[noitemsep, nosep]
						\item[1.3.] Actualiza la regla.
						\item[1.4.] Actualiza \IN{regla}.
					\end{enumerate} \\
					\hline
				\end{tabular} 
			\end{center}
		\end{table}
	
	\begin{table}[H]
		\begin{center}
			\begin{tabular}{|p{0.45\textwidth}|p{0.45\textwidth}|}
				\hline
				\multicolumn{2}{|c|}{\textbf{Flujo alterno 2: Inválido}} \\ \hline
				\textbf{Actor} & \textbf{Sistema} \\
				\hline
				\begin{enumerate}[noitemsep, nosep]
					\item[2.1.] Ingresa la regla fuera del rango de reglas en el campo \IN{regla} en \UI{configuracion_ui}.
					\item[2.2.] Presiona el botón \BTN{def_regla} mientras la simulación está activa.
				\end{enumerate} 
				& 
				\begin{enumerate}[noitemsep, nosep]
					\item[2.3.] La regla no es compatible, se ajusta a la configuración del CA.
					\item[2.4.] Actualiza la regla.
				\end{enumerate} \\
				\hline
			\end{tabular} 
		\end{center}
	\end{table}

\subsubsection*{UC17 Cambiar color de estado}

\begin{table}[H]
	\begin{center}
		\begin{tabular}{| p{0.2\textwidth} | p{0.7\textwidth} |}
			\hline 
			\textbf{Actor} & Usuario. \\ \hline
			\textbf{Descripción} & Se define un color asociado a cada estado.  \\ \hline
			\textbf{Tipo} & Secundario. \\ \hline
			\textbf{Entradas} &
			\begin{itemize}[noitemsep, nosep]
				\item Desde mouse \BTN{estado}.
			\end{itemize}
			\\ \hline
			\textbf{Salidas} &  
			\begin{itemize}[noitemsep, nosep]
				\item Interfaz \UI{simulacion_ui} actualizada.
			\end{itemize} \\ \hline
			\textbf{Precondición} & 
			\begin{itemize}[noitemsep, nosep]
				\item \BTN{estado} definido por los estados.
			\end{itemize} \\ \hline
			\textbf{Postcondición} &  
			\begin{itemize}[noitemsep, nosep]
				\item Se le asigna un color al estado seleccionado.
				\item Todas las células con ese estado actualizan su color en \UI{simulacion_ui}.
				\item La configuración de color queda almacenada.
			\end{itemize}  \\ \hline
			\textbf{Errores} & Ninguno. \\ \hline
		\end{tabular}
	\end{center}
\end{table}

\begin{table}[H]
	\begin{center}
		\begin{tabular}{|p{0.45\textwidth}|p{0.45\textwidth}|}
			\hline
			\multicolumn{2}{|c|}{\textbf{Flujo básico}} \\ \hline
			\textbf{Actor} & \textbf{Sistema} \\
			\hline
			\begin{enumerate}[noitemsep, nosep]
				\item Presiona dos veces un \BTN{estado} en \UI{configuracion_ui}.
				\item[2. ] Selecciona un color.
			\end{enumerate} 
			& 
			\begin{enumerate}[noitemsep, nosep]
				\item[3.] Abre un panel de colores. 
				\item[4.] Asigna ese color a ese estado.
				\item[5.] Cambia el color de todas las células con ese estado en \UI{simulacion_ui}.
			\end{enumerate} \\
			\hline
		\end{tabular} \\
	\end{center}
\end{table}


\subsection{Interacción con el espacio celular}

El usuario puede interactuar con el espacio celular para cambiar valores, colores o vecindades. Los siguientes casos de uso están diseñados para mejorar la visibilidad del espacio celular.

\subsubsection*{UC18 Cambiar estado de una célula}
\begin{table}[H]
	\begin{center}
		\begin{tabular}{| p{0.2\textwidth} | p{0.7\textwidth} |}
			\hline 
			\textbf{Actor} & Usuario. \\ \hline
			\textbf{Descripción} & Se cambia manualmente la configuración del espacio celular, seleccionando el estado de cada célula.  \\ \hline
			\textbf{Tipo} & Secundario. \\ \hline
			\textbf{Entradas} & \begin{enumerate}[noitemsep, nosep]
				\item Desde mouse \BTN{estado}.
				\item Desde mouse \IN{espacio}.
			\end{enumerate} 
			\\ \hline
			\textbf{Salidas} & \begin{enumerate}[noitemsep, nosep]
				\item Actualización del espacio \UI{simulacion_ui}.
			\end{enumerate}
			\\ \hline
			\textbf{Precondición} & \begin{enumerate}[noitemsep, nosep]
				\item Haber seleccionado un \BTN{estado}.
			\end{enumerate} 
			\\ \hline
			\textbf{Postcondición} & \begin{enumerate}[noitemsep, nosep]
				\item La célula cambia de estado.
			\end{enumerate}  \\ \hline
			\textbf{Errores} & \begin{enumerate}[noitemsep, nosep]
				\item Sin estado: No se ha seleccionado un estado.
			\end{enumerate} 
			\\ \hline
		\end{tabular}
	\end{center}
\end{table}

\begin{table}[H]
	\begin{center}
		\begin{tabular}{|p{0.45\textwidth}|p{0.45\textwidth}|}
			\hline
			\multicolumn{2}{|c|}{\textbf{Flujo básico}} \\ \hline
			\textbf{Actor} & \textbf{Sistema} \\
			\hline
			\begin{enumerate}[noitemsep, nosep]
				\item Selecciona un estado \BTN{estado} desde \UI{configuracion_ui}.
				\item Selecciona una célula desde \UI{simulacion_ui}.
			\end{enumerate} 
			& 
			\begin{enumerate}[noitemsep, nosep]
				\item[3.] La célula se le asigna el estado seleccionado.
				\item[4.] Se actualiza \UI{simulacion_ui}.
			\end{enumerate} \\
			\hline
		\end{tabular} \\
	\end{center}
\end{table}

\begin{table}[H]
	\begin{center}
		\begin{tabular}{|p{0.45\textwidth}|p{0.45\textwidth}|}
			\hline
			\multicolumn{2}{|c|}{\textbf{Flujo alterno 1: Sin estado}} \\ \hline
			\textbf{Actor} & \textbf{Sistema} \\
			\hline
			\begin{enumerate}[noitemsep, nosep]
				\item[1.1.] Selecciona una célula desde \UI{simulacion_ui} sin haber seleccionado un \BTN{estado}.
			\end{enumerate} 
			& 
			\begin{enumerate}[noitemsep, nosep]
				\item[1.2.] La célula cambia al estado siguiente. 
				\item[1.3.] Se actualiza \UI{simulacion_ui}.
			\end{enumerate} \\
			\hline
		\end{tabular} \\
	\end{center}
\end{table}


\subsection{Entrenamiento}

Los casos de uso siguientes permiten entrenar y definir la red neuronal del autómata celular neuronal. Se realiza el entrenamiento y la exportación de los pesos neuronales. 

\subsubsection*{UC19 Iniciar entrenamiento}

\begin{table}[H]
	\begin{center}
		\begin{tabular}{| p{0.2\textwidth} | p{0.7\textwidth} |}
			\hline 
			\textbf{Actor} & Usuario. \\ \hline
			\textbf{Descripción} & El usuario carga la evolución, define la vecindad, la configuración de la red neuronal, los hiperparámetros de entrenamiento y la ruta de guardado.  \\ \hline
			\textbf{Tipo} & Primario. \\ \hline
			\textbf{Entradas} & \begin{enumerate}[noitemsep, nosep]
				\item Desde mouse \BTN{cargar_evolucion}
				\item Desde mouse \BTN{definir_guardado}
				
				\item Desde mouse \IN{entrenamiento_vecindad}
				
				\item Desde teclado \IN{capas_ocultas}
				\item Desde teclado \IN{neuronas}
				\item Desde teclado \IN{relu}
				\item Desde teclado \IN{sigmoid}
				\item Desde teclado \IN{tanh}
				\item Desde teclado \IN{lineal}
				\item Desde teclado \IN{salida}
				\item Desde teclado \IN{funcion_salida}
				\item Desde teclado \IN{epocas}
				\item Desde teclado \IN{tamaño_lote}
				\item Desde teclado \IN{tasa_aprendizaje}
				\item Desde teclado \IN{perdida_objetivo}
				\item Desde teclado \IN{imprimir}
				\item Desde mouse \BTN{entrenar}
			\end{enumerate} 
			\\ \hline
			\textbf{Salidas} & \begin{enumerate}[noitemsep, nosep]
				\item \OUT{ruta_in}
				\item \OUT{ruta_out}
				\item \OUT{entrenamiento_vecindad_mapa}
				\item \OUT{consola}
			\end{enumerate} 
			\\ \hline
			\textbf{Precondición} & Ninguna.
			\\ \hline
			\textbf{Postcondición} & Ninguna.  \\ \hline
			\textbf{Errores} & \begin{enumerate}[noitemsep, nosep]
				\item Formato inválido: El archivo de evolución no coincide.
				\item Entrenamiento fallido: El entrenamiento no se completó debido a un error de ejecución.
			\end{enumerate} 
			\\ \hline
		\end{tabular}
	\end{center}
\end{table}

\begin{table}[H]
	\begin{center}
		\begin{tabular}{|p{0.45\textwidth}|p{0.45\textwidth}|}
			\hline
			\multicolumn{2}{|c|}{\textbf{Flujo básico}} \\ \hline
			\textbf{Actor} & \textbf{Sistema} \\
			\hline
			\begin{enumerate}[noitemsep, nosep]
				\item Presiona el botón \BTN{cargar_evolucion} desde \UI{training_ui}.
				\item[3.] Selecciona la evolución en el archivo ``npy''.
				\item[5.] Presiona el botón \BTN{definir_guardado}.
				\item[6.] Selecciona la ruta de guardado de los pesos ``npz''. 
				\item[8.] Selecciona la vecindad con \IN{entrenamiento_vecindad}.
				\item[9.] Define la configuración de la red neuronal con las siguientes entradas: \IN{capas_ocultas}, \IN{neuronas}, \IN{relu}, \IN{sigmoid}, \IN{tanh}, \IN{lineal}, \IN{salida} y \IN{funcion_salida}.
				\item[10.] Define los hiperparámetros de entrenamiento: \IN{epocas}, \IN{tamaño_lote}, \IN{tasa_aprendizaje}, \IN{perdida_objetivo}, \IN{imprimir}.
				\item[11.] Presiona el botón \BTN{entrenar}.
			\end{enumerate} 
			& 
			\begin{enumerate}[noitemsep, nosep]
				\item[2.] Abre el explorador de archivos.
				\item[4.] Muestra la ruta en \OUT{ruta_in}.
				\item[5.] Abre el explorador de archivos.
				\item[7.] Muestra la ruta en \OUT{ruta_out}.
				\item[12.] Borra el contenido de \OUT{consola}.
				\item[13.] Llama la función de entrenamiento con todos los parámetros definidos y muestra los resultados en \OUT{consola}.
				\item[14.] Guarda la evolución en la ruta definida en \OUT{ruta_out}.
			\end{enumerate} \\
			\hline
		\end{tabular} \\
	\end{center}
\end{table}

\begin{table}[H]
	\begin{center}
		\begin{tabular}{|p{0.45\textwidth}|p{0.45\textwidth}|}
			\hline
			\multicolumn{2}{|c|}{\textbf{Flujo alterno 1: Formato inválido}} \\ \hline
			\textbf{Actor} & \textbf{Sistema} \\
			\hline
			\begin{enumerate}[noitemsep, nosep]
				\item[1.1.] Presiona el botón \BTN{cargar_evolucion}.
				\item[1.3.] Selecciona un archivo incorrecto o con otro formato.
				\item[1.5.] Presiona el botón \BTN{definir_guardado}.
				\item[1.6.] Selecciona la ruta de guardado de los pesos ``npz''. 
			\end{enumerate} 
			& 
			\begin{enumerate}[noitemsep, nosep]
				\item[1.2.] Abre el explorador de archivos.
				\item[1.4.] Muestra el mensaje \MSG{formato_incorrecto}.
				\end{enumerate} \\
			\hline
		\end{tabular} \\
	\end{center}
\end{table}

\begin{table}[H]
	\begin{center}
		\begin{tabular}{|p{0.45\textwidth}|p{0.45\textwidth}|}
			\hline
			\multicolumn{2}{|c|}{\textbf{Flujo alterno 2: Entrenamiento fallido}} \\ \hline
			\textbf{Actor} & \textbf{Sistema} \\
			\hline
			\begin{enumerate}[noitemsep, nosep]
				\item[2.1.] Presiona el botón \BTN{cargar_evolucion}.
				\item[2.3.] Selecciona la evolución en el archivo ``npy''.
				\item[2.5.] Presiona el botón \BTN{definir_guardado}.
				\item[2.6.] Selecciona la ruta de guardado de los pesos ``npz''. 
				\item[2.8.] Selecciona la vecindad con \IN{entrenamiento_vecindad}.
				\item[2.9.] Define la configuración de la red neuronal con las siguientes entradas: \IN{capas_ocultas}, \IN{neuronas}, \IN{relu}, \IN{sigmoid}, \IN{tanh}, \IN{lineal}, \IN{salida} y \IN{funcion_salida}.
				\item[2.10.] Define los hiperparámetros de entrenamiento: \IN{epocas}, \IN{tamaño_lote}, \IN{tasa_aprendizaje}, \IN{perdida_objetivo}, \IN{imprimir}.
				\item[2.11.] Presiona el botón \BTN{entrenar}.
			\end{enumerate} 
			& 
			\begin{enumerate}[noitemsep, nosep]
				\item[2.2.] Abre el explorador de archivos.
				\item[2.4.] Muestra la ruta en \OUT{ruta_in}.
				\item[2.5.] Abre el explorador de archivos.
				\item[2.7.] Muestra la ruta en \OUT{ruta_out}.
				\item[2.12.] Borra el contenido de \OUT{consola}.
				\item[2.13.] Llama la función de entrenamiento con todos los parámetros definidos y muestra los resultados en \OUT{consola}.
				\item[2.14.] Algún elemento de la configuración, la cantidad de datos usados o limitaciones de hardware impidieron que se completara el entrenamiento. Muestra el mensaje \MSG{entrenamiento_incompleto} en \OUT{consola} y no guarda nada. 
			\end{enumerate} \\
			\hline
		\end{tabular} \\
	\end{center}
\end{table}

\subsubsection*{UC20 Detener entrenamiento}

\begin{table}[H]
	\begin{center}
		\begin{tabular}{| p{0.2\textwidth} | p{0.7\textwidth} |}
			\hline 
			\textbf{Actor} & Usuario. \\ \hline
			\textbf{Descripción} & El usuario detiene el entrenamiento de la red neuronal.  \\ \hline
			\textbf{Tipo} & Secundario. \\ \hline
			\textbf{Entradas} & \begin{enumerate}[noitemsep, nosep]
				\item Desde mouse \BTN{detener}.
			\end{enumerate} 
			\\ \hline
			\textbf{Salidas} & \begin{enumerate}[noitemsep, nosep]
				\item \OUT{consola}.
			\end{enumerate} 
			\\ \hline
			\textbf{Precondición} & \begin{enumerate}[noitemsep, nosep]
				\item UC18 Iniciar entrenamiento.
			\end{enumerate} 
			\\ \hline
			\textbf{Postcondición} & Ninguna.  \\ \hline
			\textbf{Errores} & Ninguno. 
			\\ \hline
		\end{tabular}
	\end{center}
\end{table}

\begin{table}[H]
	\begin{center}
		\begin{tabular}{|p{0.45\textwidth}|p{0.45\textwidth}|}
			\hline
			\multicolumn{2}{|c|}{\textbf{Flujo básico}} \\ \hline
			\textbf{Actor} & \textbf{Sistema} \\
			\hline
			\begin{enumerate}[noitemsep, nosep]
				\item Presiona el botón \BTN{detener} desde \UI{training_ui}.
			\end{enumerate} 
			& 
			\begin{enumerate}[noitemsep, nosep]
				\item[2.] Interrumpe el entrenamiento en curso.
				\item[3.] Muestra el mensaje \MSG{entrenamiento_detenido} en \OUT{consola}.
				\item[4.] Habilita nuevamente los controles de \UI{training_ui}.
			\end{enumerate} \\
			\hline
		\end{tabular}
	\end{center}
\end{table}

\subsubsection*{UC21 Borrar registro}

\begin{table}[H]
	\begin{center}
		\begin{tabular}{| p{0.2\textwidth} | p{0.7\textwidth} |}
			\hline 
			\textbf{Actor} & Usuario. \\ \hline
			\textbf{Descripción} & El usuario elimina la salida de la consola.  \\ \hline
			\textbf{Tipo} & Secundario. \\ \hline
			\textbf{Entradas} & \begin{enumerate}[noitemsep, nosep]
				\item Desde mouse \BTN{borrar_consola}.
			\end{enumerate} 
			\\ \hline
			\textbf{Salidas} & Ninguna.
			\\ \hline
			\textbf{Precondición} & Ninguna.
			\\ \hline
			\textbf{Postcondición} & Ninguna.  \\ \hline
			\textbf{Errores} & Ninguno. 
			\\ \hline
		\end{tabular}
	\end{center}
\end{table}

\begin{table}[H]
	\begin{center}
		\begin{tabular}{|p{0.45\textwidth}|p{0.45\textwidth}|}
			\hline
			\multicolumn{2}{|c|}{\textbf{Flujo básico}} \\ \hline
			\textbf{Actor} & \textbf{Sistema} \\
			\hline
			\begin{enumerate}[noitemsep, nosep]
				\item Presiona el botón \BTN{borrar_consola} desde \UI{training_ui}.
			\end{enumerate} 
			& 
			\begin{enumerate}[noitemsep, nosep]
				\item[2.] Elimina los mensajes de \OUT{consola}.
			\end{enumerate} \\
			\hline
		\end{tabular}
	\end{center}
\end{table}